% ============================================================
%  C: Como Programar -- 6ª Edição | Deitel & Deitel
%  Capítulo 1 -- Introdução aos Computadores, à Internet e à Web
%  EXERCÍCIOS DE AUTORREVISÃO (com respostas comentadas)
%  Estilo didático e incremental -- no espírito de Paul Deitel
% ============================================================
\documentclass[12pt,a4paper]{article}

% ---------- Codificação e idioma ----------
\usepackage[utf8]{inputenc}
\usepackage[T1]{fontenc}
\usepackage[brazil]{babel}

% ---------- Layout ----------
\usepackage{geometry}
\geometry{top=2.5cm, bottom=2.5cm, left=3cm, right=3cm}
\usepackage{parskip}
\setlength{\parskip}{6pt}
\setlength{\parindent}{0pt}

% ---------- Tipografia ----------



% ---------- Cores e caixas ----------
\usepackage{xcolor}
\usepackage{tcolorbox}
\tcbuselibrary{skins, breakable, theorems}

% Paleta de cores Deitel
\definecolor{deitelblue}{RGB}{0, 70, 127}
\definecolor{deitelgray}{RGB}{240, 242, 245}
\definecolor{deitelgreen}{RGB}{0, 110, 60}
\definecolor{deitellorange}{RGB}{200, 90, 0}
\definecolor{deitellightblue}{RGB}{220, 235, 250}
\definecolor{deitelred}{RGB}{160, 20, 20}

% ---------- Caixas de destaque ----------
\newtcolorbox{boaprática}[1][]{
  colback=deitellightblue, colframe=deitelblue,
  fonttitle=\bfseries\small, title={Boa prática de programação},
  breakable, #1
}
\newtcolorbox{erroco}[1][]{
  colback=red!5, colframe=deitelred,
  fonttitle=\bfseries\small, title={Erro comum de programação},
  breakable, #1
}
\newtcolorbox{portab}[1][]{
  colback=green!5, colframe=deitelgreen,
  fonttitle=\bfseries\small, title={Dica de portabilidade},
  breakable, #1
}
\newtcolorbox{respostabox}[1][]{
  colback=deitelgray, colframe=deitelblue!60,
  fonttitle=\bfseries\small\color{deitelblue}, title={Resposta detalhada},
  breakable, #1
}
\newtcolorbox{enunciadoBox}[1][]{
  colback=white, colframe=deitelblue,
  fonttitle=\bfseries, title={#1},
  breakable
}

% ---------- Listas ----------
\usepackage{enumitem}
\setlist[enumerate,1]{label=\textbf{\alph*)}, leftmargin=2em}
\setlist[itemize]{leftmargin=2em}

% ---------- Código ----------
\usepackage{listings}
\lstset{
  basicstyle=\ttfamily\small,
  backgroundcolor=\color{deitelgray},
  frame=single, framerule=0.4pt,
  rulecolor=\color{deitelblue!40},
  breaklines=true,
  showstringspaces=false,
  keywordstyle=\color{deitelblue}\bfseries,
  commentstyle=\color{deitelgreen}\itshape,
  stringstyle=\color{deitelred},
  language=C,
  numbers=none,
  xleftmargin=10pt, xrightmargin=5pt,
  belowskip=4pt, aboveskip=4pt,
}

% ---------- Cabeçalho / rodapé ----------
\usepackage{fancyhdr}
\pagestyle{fancy}
\fancyhf{}
\fancyhead[L]{\small\color{deitelblue}\textbf{C: Como Programar -- 6ª ed.\ $|$ Deitel \& Deitel}}
\fancyhead[R]{\small\color{deitelblue}Cap.\ 1 -- Autorrevisão}
\fancyfoot[C]{\small\thepage}
\renewcommand{\headrulewidth}{0.4pt}

% ---------- Títulos ----------
\usepackage{titlesec}
\titleformat{\section}[block]
  {\large\bfseries\color{deitelblue}}
  {\thesection.}{0.5em}{}[\titlerule]
\titleformat{\subsection}[block]
  {\normalsize\bfseries\color{deitelblue!80}}
  {\thesubsection.}{0.5em}{}

% ---------- Hiperlinks ----------
\usepackage{hyperref}
\hypersetup{colorlinks=true, linkcolor=deitelblue, urlcolor=deitelblue}

% ============================================================
\begin{document}

% ---- Capa ----
\begin{titlepage}
  \centering
  \vspace*{2cm}
  {\color{deitelblue}\rule{\linewidth}{2pt}}\\[0.4cm]
  {\huge\bfseries\color{deitelblue} C: Como Programar}\\[0.2cm]
  {\large\color{deitelblue} 6ª Edição $\cdot$ Paul Deitel \& Harvey Deitel}\\[0.2cm]
  {\color{deitelblue}\rule{\linewidth}{2pt}}\\[1cm]
  {\LARGE\bfseries Capítulo 1}\\[0.3cm]
  {\Large\color{deitelblue!80} Introdução aos Computadores, à Internet e à Web}\\[0.8cm]
  \begin{tcolorbox}[colback=deitellightblue, colframe=deitelblue, width=0.8\linewidth]
    \centering
    {\large\bfseries Exercícios de Autorrevisão}\\[0.2cm]
    {\normalsize Respostas completamente explicadas, passo a passo,\\
    no estilo incremental Deitel}
  \end{tcolorbox}
  \vfill
  {\small Conteúdo restrito ao Capítulo 1 $\cdot$ Abordagem incremental}
\end{titlepage}

% ---- Introdução metodológica ----
\section*{Nota ao Estudante}

Bem-vindo ao primeiro conjunto de exercícios de autorrevisão do livro \textit{C: Como Programar}.
Antes de prosseguir, queremos ressaltar a importância desta seção. Os exercícios de autorrevisão
são concebidos para que você possa verificar, de forma imediata e autônoma, se os conceitos
fundamentais do capítulo foram assimilados. Cada questão é respondida de maneira detalhada,
com explicações que vão além do simples gabarito, sempre conectando a resposta ao conteúdo
desenvolvido no Capítulo~1.

\begin{boaprática}
  Leia cada enunciado com atenção antes de consultar a resposta. Tente
  responder de memória. Essa prática ativa de recuperação de conteúdo é uma
  das técnicas mais eficazes de aprendizado comprovadas pela ciência cognitiva.
\end{boaprática}

\newpage
\tableofcontents
\newpage

% =============================================================
\section{Exercício 1.1 -- Preenchimento de Lacunas: Conceitos Gerais}
% =============================================================

\begin{enunciadoBox}[Enunciado]
Preencha os espaços em cada uma das seguintes sentenças:
\begin{enumerate}[label=\textbf{\alph*)}]
  \item A empresa que popularizou a computação pessoal foi \underline{\hspace{3cm}}.
  \item O computador que tornou a computação pessoal possível no comércio e na indústria foi o \underline{\hspace{3cm}}.
  \item Computadores processam dados sob o controle de conjuntos de instruções chamadas \underline{\hspace{3cm}} de computador.
  \item As seis unidades lógicas do computador são \underline{\hspace{2cm}}, \underline{\hspace{2cm}}, \underline{\hspace{2cm}}, \underline{\hspace{2cm}}, \underline{\hspace{2cm}} e \underline{\hspace{2cm}}.
  \item Os três tipos de linguagem que discutimos são \underline{\hspace{2cm}}, \underline{\hspace{2cm}} e \underline{\hspace{2cm}}.
  \item Os programas que traduzem programas em linguagem de alto nível na linguagem de máquina são chamados de \underline{\hspace{3cm}}.
  \item C é bastante conhecida como a linguagem de desenvolvimento do sistema operacional \underline{\hspace{3cm}}.
  \item O Department of Defense desenvolveu a linguagem Ada com uma capacidade chamada de \underline{\hspace{3cm}}, que permite que os programadores especifiquem as atividades que podem ocorrer em paralelo.
\end{enumerate}
\end{enunciadoBox}

\bigskip

% ---- a) ----
\subsection*{Item (a) -- A empresa que popularizou a computação pessoal}

\begin{respostabox}
\textbf{Resposta:} \textcolor{deitelblue}{\textbf{Apple}}

\medskip
\textbf{Explicação passo a passo:}

Na Seção~1.4 do Capítulo~1 aprendemos que, em 1977, a \textbf{Apple Computer} popularizou a
computação pessoal. Antes disso, os computadores eram equipamentos de grande porte, caros e
restritos a organizações governamentais, acadêmicas e grandes empresas. A Apple lançou
computadores pessoais acessíveis, abrindo as portas para que indivíduos pudessem usar
computadores em suas casas e pequenos negócios.

\medskip
É importante não confundir ``popularizar'' com ``introduzir''. A IBM introduziu o
\textit{IBM Personal Computer} (PC) em 1981 (item b), mas quem \textit{popularizou} -- isto é,
tornou a computação pessoal amplamente acessível e desejada -- foi a Apple, quatro anos antes.
\end{respostabox}

\bigskip

% ---- b) ----
\subsection*{Item (b) -- O computador que levou a computação pessoal ao comércio e à indústria}

\begin{respostabox}
\textbf{Resposta:} \textcolor{deitelblue}{\textbf{IBM Personal Computer (IBM PC)}}

\medskip
\textbf{Explicação passo a passo:}

Em 1981, a IBM -- na época o maior vendedor de computadores do mundo -- lançou o
\textbf{IBM Personal Computer}. Esse evento foi decisivo: pelo fato de a IBM já estar
consolidada em empresas, governos e indústrias (por meio de seus mainframes), a chegada do
IBM PC deu credibilidade comercial à computação pessoal. Empresas passaram a adotar
computadores pessoais em larga escala, transformando para sempre os ambientes de trabalho.

\medskip
Note a distinção histórica: a Apple democratizou o uso \textit{pessoal}; o IBM PC
democratizou o uso \textit{corporativo}. Ambos, juntos, definiram a era da computação pessoal.
\end{respostabox}

\bigskip

% ---- c) ----
\subsection*{Item (c) -- Conjuntos de instruções que controlam os computadores}

\begin{respostabox}
\textbf{Resposta:} \textcolor{deitelblue}{\textbf{programas}}

\medskip
\textbf{Explicação passo a passo:}

Da Seção~1.2: \textit{``Computadores processam dados sob o controle de conjuntos de
instruções chamados \textbf{programas} de computador.''}

Essa definição é fundamental. Um computador, por si só, é apenas hardware -- circuitos,
chips e dispositivos físicos. O que dá vida a esse hardware e o faz executar tarefas
úteis são os \textbf{programas} (software) escritos por programadores. Programar em C é
exatamente a arte de criar esses conjuntos de instruções ordenadas.
\end{respostabox}

\bigskip

% ---- d) ----
\subsection*{Item (d) -- As seis unidades lógicas do computador}

\begin{respostabox}
\textbf{Resposta:} \textcolor{deitelblue}{\textbf{(1) unidade de entrada, (2) unidade de saída, (3) unidade de memória, (4) unidade lógica e aritmética (ALU), (5) unidade central de processamento (CPU), (6) unidade de armazenamento secundário}}

\medskip
\textbf{Explicação passo a passo:}

A Seção~1.3 descreve que, independentemente de sua aparência física, todo computador pode
ser dividido em \textbf{seis unidades lógicas}. Veja cada uma delas:

\begin{itemize}
  \item \textbf{Unidade de entrada:} A seção ``receptora'' do computador. Recebe dados e
  programas do mundo externo (teclado, mouse, sensores, etc.) e os disponibiliza para as
  demais unidades.

  \item \textbf{Unidade de saída:} A seção de ``expedição''. Pega as informações processadas
  e as envia para dispositivos externos (monitor, impressora, alto-falantes, etc.).

  \item \textbf{Unidade de memória:} O ``depósito rápido'' do computador. Armazena dados e
  programas temporariamente enquanto estão em uso. Suas informações são \textit{voláteis}
  (perdidas quando o computador é desligado). Também chamada de \textit{memória primária}
  ou \textit{RAM}.

  \item \textbf{Unidade lógica e aritmética (ALU):} A seção de ``processamento''. Realiza
  cálculos matemáticos (adição, subtração, multiplicação, divisão) e tomada de decisões
  lógicas (comparações entre valores).

  \item \textbf{Unidade central de processamento (CPU):} A seção ``administrativa''.
  Coordena e supervisiona todas as demais unidades. É ela que diz \textit{quando} a
  unidade de entrada deve ler dados, \textit{quando} a ALU deve calcular, \textit{quando}
  a unidade de saída deve enviar resultados.

  \item \textbf{Unidade de armazenamento secundário:} O ``depósito de longo prazo''.
  Armazena programas e dados de forma \textit{persistente} (mantida mesmo com o computador
  desligado), em discos rígidos, CDs, DVDs, pendrives etc.
\end{itemize}

\medskip
\begin{boaprática}
  Memorize as seis unidades lógicas e suas funções. Esse modelo conceitual
  é independente do hardware físico e será útil durante todo o curso, pois
  entender \textit{onde} os dados vivem (entrada → memória → ALU/CPU → saída/disco)
  é essencial para escrever programas eficientes em C.
\end{boaprática}
\end{respostabox}

\bigskip

% ---- e) ----
\subsection*{Item (e) -- Os três tipos de linguagem de programação}

\begin{respostabox}
\textbf{Resposta:} \textcolor{deitelblue}{\textbf{linguagens de máquina, linguagens simbólicas (assembly) e linguagens de alto nível}}

\medskip
\textbf{Explicação passo a passo:}

A Seção~1.6 apresenta a evolução histórica das linguagens de programação em três gerações:

\begin{enumerate}
  \item \textbf{Linguagens de máquina:} O ``idioma nativo'' do hardware. Consistem em
  sequências de números binários (0s e 1s) que a CPU entende diretamente. São
  dependentes de máquina -- um programa escrito para um tipo de CPU não funciona
  em outro. Extremamente difíceis para seres humanos lerem e escreverem.

  \item \textbf{Linguagens simbólicas (assembly):} Substituem as sequências numéricas por
  abreviações mnemônicas (como \texttt{load}, \texttt{add}, \texttt{store}). Programas
  chamados \textit{assemblers} (montadores) traduzem essas abreviações para linguagem de
  máquina. Ainda dependentes de máquina, porém muito mais legíveis que a linguagem de máquina pura.

  \item \textbf{Linguagens de alto nível:} Usam palavras próximas ao inglês e notação
  matemática convencional (como \texttt{valorTotal = salario + horaExtra;}). Uma única
  instrução de alto nível pode gerar dezenas de instruções de máquina. Programas chamados
  \textit{compiladores} fazem essa tradução. \textbf{C é uma linguagem de alto nível.}
\end{enumerate}
\end{respostabox}

\bigskip

% ---- f) ----
\subsection*{Item (f) -- Programas que traduzem linguagem de alto nível para linguagem de máquina}

\begin{respostabox}
\textbf{Resposta:} \textcolor{deitelblue}{\textbf{compiladores}}

\medskip
\textbf{Explicação passo a passo:}

Da Seção~1.6: programas tradutores chamados \textbf{compiladores} convertem os programas
escritos em linguagem de alto nível (como C) para a linguagem de máquina que o computador
pode executar.

É importante distinguir dois tipos de tradutores:
\begin{itemize}
  \item \textbf{Compiladores:} Traduzem o programa \textit{inteiro} de uma vez, gerando um
  arquivo executável (código objeto). A execução posterior é rápida.
  \item \textbf{Interpretadores:} Traduzem e executam o programa \textit{linha por linha},
  sem gerar um executável separado. São mais lentos na execução, mas eliminam a etapa
  de compilação.
\end{itemize}

Em C, usamos \textit{compiladores} (como o \texttt{gcc} no Linux ou o Visual C++ no Windows).
\end{respostabox}

\bigskip

% ---- g) ----
\subsection*{Item (g) -- O sistema operacional associado ao desenvolvimento de C}

\begin{respostabox}
\textbf{Resposta:} \textcolor{deitelblue}{\textbf{UNIX}}

\medskip
\textbf{Explicação passo a passo:}

Da Seção~1.7: C foi criada por \textbf{Dennis Ritchie} no Bell Laboratories, em 1972,
originalmente implementada em um computador DEC PDP-11. Sua principal motivação foi o
desenvolvimento do \textbf{sistema operacional UNIX}. C e UNIX cresceram juntos: C tornou-se
a linguagem de implementação do UNIX e ganhou popularidade com ele.

Hoje, praticamente todos os sistemas operacionais modernos -- Linux, macOS, iOS, Android --
têm raízes no UNIX e foram escritos em C e/ou C++. Isso ilustra o alcance e a longevidade
de C como linguagem de sistemas.
\end{respostabox}

\bigskip

% ---- h) ----
\subsection*{Item (h) -- O recurso da linguagem Ada para paralelismo}

\begin{respostabox}
\textbf{Resposta:} \textcolor{deitelblue}{\textbf{multitarefa} (do inglês \textit{multitasking})}

\medskip
\textbf{Explicação passo a passo:}

Da Seção~1.11: a linguagem \textbf{Ada} foi desenvolvida pelo Departamento de Defesa dos
EUA (DoD) e possui um recurso chamado \textbf{multitarefa}, que permite que os programadores
especifiquem que muitas atividades devem ocorrer simultaneamente (em paralelo).

A Ada recebeu seu nome em homenagem a \textbf{Lady Ada Lovelace}, filha do poeta Lord Byron,
reconhecida como a primeira programadora de computadores da história. Note que, embora a
multitarefa não faça parte do C padrão (ANSI C), ela está disponível por meio de bibliotecas
complementares, como a POSIX Threads (\texttt{pthreads}).
\end{respostabox}

\newpage

% =============================================================
\section{Exercício 1.2 -- Preenchimento de Lacunas: Ambiente de Desenvolvimento em C}
% =============================================================

\begin{enunciadoBox}[Enunciado]
Preencha os espaços em cada uma das sentenças a respeito do ambiente C:
\begin{enumerate}[label=\textbf{\alph*)}]
  \item Programas em C são normalmente digitados em um computador utilizando-se um programa \underline{\hspace{3cm}}.
  \item Em um sistema C, um programa \underline{\hspace{3cm}} é executado automaticamente antes que se inicie a fase de tradução.
  \item Os tipos mais comuns de diretivas de pré-processador são utilizados para \underline{\hspace{3cm}} e \underline{\hspace{3cm}}.
  \item O programa \underline{\hspace{3cm}} combina a saída do compilador com várias funções de biblioteca para produzir uma imagem executável.
  \item O programa \underline{\hspace{3cm}} transfere a imagem executável do disco para a memória.
  \item Para carregar e executar o programa compilado mais recentemente em um sistema Linux, digite \underline{\hspace{3cm}}.
\end{enumerate}
\end{enunciadoBox}

\bigskip

% ---- a) ----
\subsection*{Item (a) -- O programa usado para digitar código C}

\begin{respostabox}
\textbf{Resposta:} \textcolor{deitelblue}{\textbf{editor} (programa editor)}

\medskip
\textbf{Explicação passo a passo:}

A \textbf{Fase 1} do desenvolvimento de um programa em C, descrita na Seção~1.14
(Figura~1.1), consiste em \textit{editar} o arquivo fonte. Um \textbf{programa editor}
é utilizado para criar e modificar o código-fonte.

Exemplos de editores amplamente usados:
\begin{itemize}
  \item No Linux: \texttt{vi}, \texttt{vim}, \texttt{emacs}, \texttt{nano}, \texttt{gedit}
  \item Em ambientes integrados (IDEs): o editor já vem embutido (ex.: Microsoft Visual
  Studio, Eclipse, Code::Blocks)
\end{itemize}

Os arquivos de código-fonte C devem ter extensão \textbf{\texttt{.c}} (por exemplo,
\texttt{bemvindo.c}, \texttt{calculadora.c}).
\end{respostabox}

\bigskip

% ---- b) ----
\subsection*{Item (b) -- O programa executado automaticamente antes da compilação}

\begin{respostabox}
\textbf{Resposta:} \textcolor{deitelblue}{\textbf{pré-processador}}

\medskip
\textbf{Explicação passo a passo:}

Na \textbf{Fase 2} (descrita na Seção~1.14), antes que o compilador propriamente dito
inicie a tradução do código C para linguagem de máquina, um programa especial chamado
\textbf{pré-processador C} é executado \textit{automaticamente}.

O pré-processador obedece a comandos especiais chamados \textbf{diretivas de
pré-processador}, que começam com o caractere \texttt{\#}. O pré-processador prepara
o código-fonte para a compilação, realizando as manipulações indicadas por essas
diretivas.

\begin{boaprática}
  Em um sistema C típico, o programador não precisa invocar o pré-processador
  manualmente. Ao executar o comando de compilação (\texttt{cc} ou \texttt{gcc}),
  o pré-processador é acionado automaticamente como primeira etapa.
\end{boaprática}
\end{respostabox}

\bigskip

% ---- c) ----
\subsection*{Item (c) -- Os usos mais comuns das diretivas de pré-processador}

\begin{respostabox}
\textbf{Resposta:} \textcolor{deitelblue}{\textbf{incluir outros arquivos no arquivo a ser compilado} e \textbf{substituir símbolos especiais por texto de programa}}

\medskip
\textbf{Explicação passo a passo:}

As \textbf{diretivas de pré-processador} mais comuns, conforme a Seção~1.14, servem para:

\begin{enumerate}
  \item \textbf{Incluir outros arquivos:} A diretiva \texttt{\#include} insere o conteúdo de
  outro arquivo no arquivo a ser compilado. Por exemplo:
  \begin{lstlisting}
#include <stdio.h>  /* inclui funcoes de entrada/saida */
  \end{lstlisting}
  Isso é essencial para usar funções da biblioteca-padrão de C, como \texttt{printf} e
  \texttt{scanf}.

  \item \textbf{Substituições de texto (macros):} A diretiva \texttt{\#define} cria
  \textit{macros} -- símbolos que serão substituídos por um valor ou trecho de código
  antes da compilação. Por exemplo:
  \begin{lstlisting}
#define PI 3.14159265
  \end{lstlisting}
  Toda ocorrência de \texttt{PI} no código será substituída por \texttt{3.14159265}
  antes que o compilador processe o arquivo.
\end{enumerate}

Uma discussão detalhada de todas as diretivas do pré-processador é apresentada no
Capítulo~13 do livro.
\end{respostabox}

\bigskip

% ---- d) ----
\subsection*{Item (d) -- O programa que combina código objeto com bibliotecas}

\begin{respostabox}
\textbf{Resposta:} \textcolor{deitelblue}{\textbf{editor de ligação} (ou \textit{linker})}

\medskip
\textbf{Explicação passo a passo:}

A \textbf{Fase 4} do processo é chamada de \textit{ligação} (\textit{linking}). O
\textbf{editor de ligação} (ou \textit{linker}) resolve as referências externas presentes
no código objeto gerado pelo compilador.

Quando o compilador traduz seu código C para código objeto, ele deixa ``lacunas'' para
as funções que você usou, mas que estão definidas em outros arquivos (por exemplo, na
biblioteca-padrão de C). O \textit{linker} preenche essas lacunas, combinando o código
objeto do seu programa com o código das funções de biblioteca, produzindo uma
\textbf{imagem executável} completa -- um arquivo que pode ser executado pelo sistema
operacional.

No Linux, tanto a compilação quanto a ligação são feitas pelo comando \texttt{gcc}:
\begin{lstlisting}
gcc bemvindo.c   /* compila E liga, gerando a.out */
\end{lstlisting}
\end{respostabox}

\bigskip

% ---- e) ----
\subsection*{Item (e) -- O programa que transfere a imagem executável para a memória}

\begin{respostabox}
\textbf{Resposta:} \textcolor{deitelblue}{\textbf{carregador} (\textit{loader})}

\medskip
\textbf{Explicação passo a passo:}

A \textbf{Fase 5} é a \textit{carga} (\textit{loading}). O \textbf{carregador}
(\textit{loader}) pega a imagem executável armazenada no disco e a transfere para a
\textbf{memória primária} (RAM), deixando-a pronta para ser executada pela CPU.

Além do executável principal, o carregador também carrega os componentes adicionais das
\textbf{bibliotecas compartilhadas} (\textit{shared libraries}) necessários ao programa.
Em sistemas modernos, o carregador é chamado automaticamente pelo sistema operacional
quando o usuário solicita a execução do programa.
\end{respostabox}

\bigskip

% ---- f) ----
\subsection*{Item (f) -- Comando para executar o programa em Linux}

\begin{respostabox}
\textbf{Resposta:} \textcolor{deitelblue}{\textbf{\texttt{./a.out}}}

\medskip
\textbf{Explicação passo a passo:}

Na \textbf{Fase 6} (execução), para carregar e executar o programa mais recentemente
compilado em um sistema Linux, digita-se no terminal:

\begin{lstlisting}
./a.out
\end{lstlisting}

Por que \texttt{./a.out}? Vamos entender cada parte:

\begin{itemize}
  \item \texttt{./} $\rightarrow$ indica ao sistema operacional que o executável está no
  \textit{diretório atual} (o ``\texttt{.}'' representa o diretório corrente e ``\texttt{/}''
  é o separador de caminho no Linux).
  \item \texttt{a.out} $\rightarrow$ é o nome padrão do arquivo executável gerado pelo
  \texttt{gcc} quando nenhum nome específico é informado.
\end{itemize}

Você pode especificar um nome diferente para o executável usando a opção \texttt{-o}:
\begin{lstlisting}
gcc bemvindo.c -o bemvindo   /* gera executavel chamado "bemvindo" */
./bemvindo                   /* executa o programa */
\end{lstlisting}

\begin{erroco}
  Os comandos do Linux diferenciam maiúsculas de minúsculas
  (\textit{case-sensitive}). Digitar \texttt{./A.OUT} ou \texttt{./A.out}
  resultará em erro, pois o arquivo gerado pelo \texttt{gcc} é \texttt{a.out}
  (letras minúsculas).
\end{erroco}
\end{respostabox}

\newpage

% =============================================================
\section*{Visão Geral: O Ciclo de Desenvolvimento de Programas em C}
% =============================================================

Para consolidar os conceitos abordados no Exercício~1.2, apresentamos abaixo um resumo
visual do ciclo completo de desenvolvimento de um programa em C (Figura~1.1 do livro):

\bigskip

\begin{tcolorbox}[colback=deitelgray, colframe=deitelblue, title={\bfseries As Seis Fases do Desenvolvimento em C}, breakable]
\begin{enumerate}
  \item[\textbf{Fase 1}] \textbf{Edição} $\longrightarrow$ O programador cria o arquivo
  \texttt{.c} com um \textbf{editor de texto}.

  \item[\textbf{Fase 2}] \textbf{Pré-processamento} $\longrightarrow$ O \textbf{pré-processador}
  executa as diretivas \texttt{\#include}, \texttt{\#define}, etc.

  \item[\textbf{Fase 3}] \textbf{Compilação} $\longrightarrow$ O \textbf{compilador}
  traduz o código C para código objeto (linguagem de máquina).

  \item[\textbf{Fase 4}] \textbf{Ligação (Linking)} $\longrightarrow$ O \textbf{editor de
  ligação} combina o código objeto com as bibliotecas, gerando a imagem executável
  (\texttt{a.out}).

  \item[\textbf{Fase 5}] \textbf{Carregamento (Loading)} $\longrightarrow$ O
  \textbf{carregador} transfere o executável do disco para a memória RAM.

  \item[\textbf{Fase 6}] \textbf{Execução} $\longrightarrow$ A \textbf{CPU} executa o
  programa, instrução por instrução.
\end{enumerate}
\end{tcolorbox}

\bigskip

\begin{portab}
  Em sistemas Linux, todo esse processo (Fases~2 a~4) é normalmente ativado por um
  único comando: \texttt{gcc nome.c}. A simplicidade do processo é uma das razões
  pelas quais C permanece como a linguagem preferida para programação de sistemas
  em múltiplas plataformas.
\end{portab}

\newpage

% =============================================================
\section*{Resumo das Respostas dos Exercícios de Autorrevisão}
% =============================================================

\begin{tcolorbox}[colback=deitellightblue, colframe=deitelblue, title={\bfseries\large Gabarito Rápido -- Capítulo 1}, breakable]

\textbf{Exercício 1.1:}
\begin{enumerate}[label=\alph*)]
  \item Apple
  \item IBM Personal Computer (IBM PC)
  \item programas
  \item unidade de entrada, unidade de saída, unidade de memória, unidade lógica e aritmética (ALU), unidade central de processamento (CPU), unidade de armazenamento secundário
  \item linguagens de máquina, linguagens simbólicas (assembly), linguagens de alto nível
  \item compiladores
  \item UNIX
  \item multitarefa
\end{enumerate}

\medskip
\textbf{Exercício 1.2:}
\begin{enumerate}[label=\alph*)]
  \item editor (programa editor)
  \item pré-processador
  \item incluir outros arquivos no arquivo a ser compilado; substituir símbolos especiais por texto de programa
  \item editor de ligação (\textit{linker})
  \item carregador (\textit{loader})
  \item \texttt{./a.out}
\end{enumerate}
\end{tcolorbox}

\bigskip

\begin{boaprática}
  Após verificar suas respostas, retorne ao Capítulo~1 e releia os trechos
  referentes aos itens que você errou. A leitura direcionada, motivada por uma
  dificuldade específica, é muito mais eficaz do que a leitura passiva e linear.
  Esse é o espírito da abordagem \textit{incremental} Deitel.
\end{boaprática}

\end{document}
